\begin{figure*}
    \centering
    \includegraphics[width=\linewidth]{figs/gpt-architecture.pdf}
    \caption{Overall dataflow architecture of a single Transformer layer that uses post-LayerNorm scheme~\cite{radford2019gpt2}.
    SDP denotes scaled dot-product.
    Orange nodes denote the GEMM kernels.
    Yellow nodes are the non-linear kernels, including softmax (SM), LayerNorm (LN), and GELU (GL).
    Green rectangles represent the FIFOs between kernels, and purple rectangles are the data loaders.}
    \label{fig:arch}
\end{figure*}

\subsection{Accelerator Design}
\label{sec:xcel}
We integrate the proposed kernels in \S\ref{subsec:kernels} to create a high-performance hardware accelerator.
The overall architecture of our proposed accelerator is depicted in Figure~\ref{fig:arch}.
Different operators including linear and non-linear operators are connected with FIFOs, except that the KV cache discussed in \S\ref{subsubsec:memory_resource} leverages double buffers.
This design reads input tensors from off-chip memory and stores the results back to memory after each layer.
Intermediate activation values are directly streamed to the next operator.
\revise{Initially, all parameters are stored in DRAM, and data loaders $L_{(\cdot)}$ are responsible for loading data from DRAM or streaming buffers.
Storing parameters off-chip reduces the number of FPGA devices needed and avoids the need to build different bitstreams for different network layers.
After completing a layer, the accelerator fetches new parameters from the host to the device for the subsequent layer.
Since data loading is hidden from the computation, overall latency remains unaffected.}
Given the contemporary trend of multi-die FPGA designs~\cite{guo2021autobridge,du2022hisparse}, explicit dataflow partitioning becomes necessary to meet timing requirements.
Our target device, Alveo U280 FPGA~\cite{u280}, has three chip dies called super logic regions (SLRs).
Consequently, we partition the dataflow into three regions to confine each region fully inside SLR.
\revise{According to our placement constraints, AMD Vitis toolchain will automatically insert AXI Register Slice IPs to pipeline SLR crossing.}

% \yx{I think there is no use to just list out what kernels are on which region. Instead, we should explain why we employ this division.}

We leverage the proposed analytical framework to guide our accelerator design.
Since typical Transformer models have $\dffn=4d$~\cite{devlin2018bert,radford2019gpt2} and $l<d$, according to work balancing of Equation~\eqref{eq:m}, we have $M_{q,k,v,p} = M$, $M_{a_1, a_2} < M$, and $M_{f1, f2} = 4M$.
A straight-forward division is to put the PEs for $q$, $k$, $v$, SDP, and $p$ on SLR0, $f1$ on SLR1, and $f2$ on SLR2 so that each SLR roughly contains $4M$ MAC units.
However, we observe that scaling up the linear operators in FFN poses significant challenges to timing closure. 
Among various configurations of systolic arrays we tested, the maximum capacity of one SLR at 250 MHz is three of 8~\X16 systolic arrays; a single 16~\X16 one fails timing. 
Therefore, we only leverage 8~\X8 and 8~\X16 systolic arrays for simplicity. 
\revise{We also explore using LUT-based multipliers as they provide greater flexibility for placement compared to DSPs.
However, the presence of additional inter-LUT wires results in a much lower frequency (191 MHz) compared to the DSP-based multipliers.}
To minimize the number of SLR crossings, we put $q$, $k$, and $v$ on SLR0 and use 8~\X16 systolic arrays, which also ensures a relatively low latency for the first stage based on Equation~\eqref{eq:latency_prefill}.
MHA and the $p$ projection are on SLR1, with $a_1$ and $a_2$ using 8~\X8, and $p$ using 8~\X16 systolic arrays.
$f_1$ and $f_2$ operators on SLR2 using 8~\X16 systolic arrays.
Therefore, it can still form a relatively balanced 3:2:2 resource utilization ratio for linear operators.
% In this case, the latency of one layer would be $5ld^2/M$ based on Equation~\eqref{eq:latency_prefill}.

% Region 0 consists of the input loader and the $q$, $k$, $v$ linear layers.
% Input data resides in the High Bandwidth Memory (HBM) next to SLR0.
% To facilitate efficient data transfer, three separate loaders read input from HBM for $k$ and $v$, $q$, and residual connection tensors.
% This division is critical because the $kv$ path exhibits slower data transfer rates compared to the other two paths, necessitating storage of all results in a double buffer for the subsequent SDP operators.
% Sharing a loader for these paths would risk deadlock scenarios, as the faster path could block the loader when the associated FIFO becomes full.
% The loader streams data to the linear layers and forwards the output to the next region.

% Region 1 includes the core SDP operators, along with the ensuing residual connection and LayerNorm operators.
% Following double buffering of the $k$ and $v$ outputs, the $k$ tensor can be streamed directly into \texttt{matmul1} for matrix multiplication.
% We calculate results for different heads sequentially to conserve memory.
% After \texttt{matmul1}, the results are streamed into the softmax function and subsequently passed to \texttt{matmul2}.
% Finally, the SDP results are concatenated and processed through the projection layer, merging with the input shortcut before undergoing the LayerNorm module.
% The resulting output is transmitted to Region 2.

% Region 2 mainly features the FFN module.
% The connection is straightforward, but careful selection of FIFO sizes is crucial. Given that the production rate of \texttt{FFN1} exceeds the consumption rate of \texttt{FFN2}, a sufficiently large FIFO is required to store intermediate results between these two layers. Inadequate buffer size would lead to buffer overflow and backpressure, causing dataflow stalls. Similarly, a sizable buffer is necessary for the shortcut to prevent deadlock. Following the residual connection and LayerNorm operator, the results are written back to memory.

% \yx{is this our true reason of doing such a division?}
% This partitioning scheme is specifically designed to maintain a balanced pipeline stage by splitting it after the LayerNorm function, facilitating efficient communication with only a single output tensor relayed from Region 1 to Region 2.
% Other partitioning points would entail communication with more than two tensors, thereby increasing communication overhead.
